% ============================================================
% shared/preamble.tex
% Zentrale Konfiguration für alle Bände der Buchreihe
% "Python für Robotik und ROS2"
% ============================================================
% Einbinden: % ============================================================
% shared/preamble.tex
% Zentrale Konfiguration für alle Bände der Buchreihe
% "Python für Robotik und ROS2"
% ============================================================
% Einbinden: % ============================================================
% shared/preamble.tex
% Zentrale Konfiguration für alle Bände der Buchreihe
% "Python für Robotik und ROS2"
% ============================================================
% Einbinden: % ============================================================
% shared/preamble.tex
% Zentrale Konfiguration für alle Bände der Buchreihe
% "Python für Robotik und ROS2"
% ============================================================
% Einbinden: \input{../shared/preamble} in band*/main.tex
% ============================================================

% ── Buchreihe-Variablen (für shared/titelseite.tex kompatibel) ────
────────────────────────────────────────
\newcommand{\BuchReihe}{Robotik entwickeln lernen}
\newcommand{\BandAutor}{Wolfgang Becker}
\newcommand{\BandVersion}{Arbeitsfassung 2026}


% ── Sprache & Schrift ────────────────────────────────────────
\usepackage[ngerman]{babel}
\usepackage[T1]{fontenc}
\usepackage{lmodern}

% ── Layout ──────────────────────────────────────────────────
\usepackage{geometry}
\geometry{left=28mm,right=22mm,top=25mm,bottom=30mm}
\usepackage{scrlayer-scrpage}
\KOMAoptions{headsepline=true}
\pagestyle{scrheadings}
\clearpairofpagestyles
\lehead{\leftmark}
\rohead{\BuchReihe}   % wird in band*/main.tex gesetzt
\cfoot{\thepage}

% ── Pakete ──────────────────────────────────────────────────
\usepackage{xcolor}
\usepackage{tcolorbox}
\tcbuselibrary{skins,breakable,listings}
\usepackage{enumitem}
\usepackage{hyperref}
\usepackage{listings}
\usepackage{amsmath}
\usepackage{amssymb}
\usepackage{textcomp}
\usepackage{graphicx}
\usepackage{booktabs}
\usepackage{tabularx}

\hypersetup{colorlinks=true,linkcolor=blue!60!black,urlcolor=teal}

% ── Farben ──────────────────────────────────────────────────
\definecolor{MerksatzFarbe}{HTML}{DDEEFF}
\definecolor{UebungFarbe}{HTML}{FFF3CC}
\definecolor{PraxisFarbe}{HTML}{E8F5E9}
\definecolor{TippFarbe}{HTML}{FFF8E1}
\definecolor{RosFarbe}{HTML}{FCE4EC}
\definecolor{AnalogieFarbe}{HTML}{EDE7F6}
\definecolor{LoesungFarbe}{HTML}{F3E5F5}
\definecolor{CodeHintergrund}{HTML}{F8F8F8}

% ── Python-Code ─────────────────────────────────────────────
\lstdefinestyle{pythonstyle}{
    language=Python,
    backgroundcolor=\color{CodeHintergrund},
    basicstyle=\ttfamily\small,
    keywordstyle=\color{blue!70!black}\bfseries,
    stringstyle=\color{teal!80!black},
    commentstyle=\color{gray!70!black}\itshape,
    showstringspaces=false,
    breaklines=true,
    frame=single,
    rulecolor=\color{gray!40},
    tabsize=4,
    extendedchars=true,
    numbers=left,
    numberstyle=\tiny\color{gray},
    numbersep=8pt
}
\lstset{
    style=pythonstyle,
    literate=
        {ä}{{\"a}}1 {ö}{{\"o}}1 {ü}{{\"u}}1
        {Ä}{{\"A}}1 {Ö}{{\"O}}1 {Ü}{{\"U}}1
        {ß}{{\ss}}1
}

% ── Didaktische Boxen ────────────────────────────────────────
\newtcolorbox{merksatz}[1][]{enhanced,breakable,
    colback=MerksatzFarbe,colframe=blue!60!black,
    fonttitle=\bfseries,title=Merksatz,#1}

\newtcolorbox{uebung}[1][]{enhanced,breakable,
    colback=UebungFarbe,colframe=orange!80!black,
    fonttitle=\bfseries,title=Übung,#1}

\newtcolorbox{praxis}[1][]{enhanced,breakable,
    colback=PraxisFarbe,colframe=green!60!black,
    fonttitle=\bfseries,title=Praxisbezug Robotik,#1}

\newtcolorbox{tipp}[1][]{enhanced,breakable,
    colback=TippFarbe,colframe=orange!40!black,
    fonttitle=\bfseries,title=Tipp,#1}

\newtcolorbox{ros}[1][]{enhanced,breakable,
    colback=RosFarbe,colframe=red!70!black,
    fonttitle=\bfseries,title=ROS2-Verbindung,#1}

\newtcolorbox{analogie}[1][]{enhanced,breakable,
    colback=AnalogieFarbe,colframe=purple!60!black,
    fonttitle=\bfseries,title=Analogie,#1}

\newtcolorbox{loesung}[1][]{enhanced,breakable,
    colback=LoesungFarbe,colframe=purple!60!black,
    fonttitle=\bfseries,title=Lösung,#1}
 in band*/main.tex
% ============================================================

% ── Buchreihe-Variablen (für shared/titelseite.tex kompatibel) ────
────────────────────────────────────────
\newcommand{\BuchReihe}{Robotik entwickeln lernen}
\newcommand{\BandAutor}{Wolfgang Becker}
\newcommand{\BandVersion}{Arbeitsfassung 2026}


% ── Sprache & Schrift ────────────────────────────────────────
\usepackage[ngerman]{babel}
\usepackage[T1]{fontenc}
\usepackage{lmodern}

% ── Layout ──────────────────────────────────────────────────
\usepackage{geometry}
\geometry{left=28mm,right=22mm,top=25mm,bottom=30mm}
\usepackage{scrlayer-scrpage}
\KOMAoptions{headsepline=true}
\pagestyle{scrheadings}
\clearpairofpagestyles
\lehead{\leftmark}
\rohead{\BuchReihe}   % wird in band*/main.tex gesetzt
\cfoot{\thepage}

% ── Pakete ──────────────────────────────────────────────────
\usepackage{xcolor}
\usepackage{tcolorbox}
\tcbuselibrary{skins,breakable,listings}
\usepackage{enumitem}
\usepackage{hyperref}
\usepackage{listings}
\usepackage{amsmath}
\usepackage{amssymb}
\usepackage{textcomp}
\usepackage{graphicx}
\usepackage{booktabs}
\usepackage{tabularx}

\hypersetup{colorlinks=true,linkcolor=blue!60!black,urlcolor=teal}

% ── Farben ──────────────────────────────────────────────────
\definecolor{MerksatzFarbe}{HTML}{DDEEFF}
\definecolor{UebungFarbe}{HTML}{FFF3CC}
\definecolor{PraxisFarbe}{HTML}{E8F5E9}
\definecolor{TippFarbe}{HTML}{FFF8E1}
\definecolor{RosFarbe}{HTML}{FCE4EC}
\definecolor{AnalogieFarbe}{HTML}{EDE7F6}
\definecolor{LoesungFarbe}{HTML}{F3E5F5}
\definecolor{CodeHintergrund}{HTML}{F8F8F8}

% ── Python-Code ─────────────────────────────────────────────
\lstdefinestyle{pythonstyle}{
    language=Python,
    backgroundcolor=\color{CodeHintergrund},
    basicstyle=\ttfamily\small,
    keywordstyle=\color{blue!70!black}\bfseries,
    stringstyle=\color{teal!80!black},
    commentstyle=\color{gray!70!black}\itshape,
    showstringspaces=false,
    breaklines=true,
    frame=single,
    rulecolor=\color{gray!40},
    tabsize=4,
    extendedchars=true,
    numbers=left,
    numberstyle=\tiny\color{gray},
    numbersep=8pt
}
\lstset{
    style=pythonstyle,
    literate=
        {ä}{{\"a}}1 {ö}{{\"o}}1 {ü}{{\"u}}1
        {Ä}{{\"A}}1 {Ö}{{\"O}}1 {Ü}{{\"U}}1
        {ß}{{\ss}}1
}

% ── Didaktische Boxen ────────────────────────────────────────
\newtcolorbox{merksatz}[1][]{enhanced,breakable,
    colback=MerksatzFarbe,colframe=blue!60!black,
    fonttitle=\bfseries,title=Merksatz,#1}

\newtcolorbox{uebung}[1][]{enhanced,breakable,
    colback=UebungFarbe,colframe=orange!80!black,
    fonttitle=\bfseries,title=Übung,#1}

\newtcolorbox{praxis}[1][]{enhanced,breakable,
    colback=PraxisFarbe,colframe=green!60!black,
    fonttitle=\bfseries,title=Praxisbezug Robotik,#1}

\newtcolorbox{tipp}[1][]{enhanced,breakable,
    colback=TippFarbe,colframe=orange!40!black,
    fonttitle=\bfseries,title=Tipp,#1}

\newtcolorbox{ros}[1][]{enhanced,breakable,
    colback=RosFarbe,colframe=red!70!black,
    fonttitle=\bfseries,title=ROS2-Verbindung,#1}

\newtcolorbox{analogie}[1][]{enhanced,breakable,
    colback=AnalogieFarbe,colframe=purple!60!black,
    fonttitle=\bfseries,title=Analogie,#1}

\newtcolorbox{loesung}[1][]{enhanced,breakable,
    colback=LoesungFarbe,colframe=purple!60!black,
    fonttitle=\bfseries,title=Lösung,#1}
 in band*/main.tex
% ============================================================

% ── Buchreihe-Variablen (für shared/titelseite.tex kompatibel) ────
────────────────────────────────────────
\newcommand{\BuchReihe}{Robotik entwickeln lernen}
\newcommand{\BandAutor}{Wolfgang Becker}
\newcommand{\BandVersion}{Arbeitsfassung 2026}


% ── Sprache & Schrift ────────────────────────────────────────
\usepackage[ngerman]{babel}
\usepackage[T1]{fontenc}
\usepackage{lmodern}

% ── Layout ──────────────────────────────────────────────────
\usepackage{geometry}
\geometry{left=28mm,right=22mm,top=25mm,bottom=30mm}
\usepackage{scrlayer-scrpage}
\KOMAoptions{headsepline=true}
\pagestyle{scrheadings}
\clearpairofpagestyles
\lehead{\leftmark}
\rohead{\BuchReihe}   % wird in band*/main.tex gesetzt
\cfoot{\thepage}

% ── Pakete ──────────────────────────────────────────────────
\usepackage{xcolor}
\usepackage{tcolorbox}
\tcbuselibrary{skins,breakable,listings}
\usepackage{enumitem}
\usepackage{hyperref}
\usepackage{listings}
\usepackage{amsmath}
\usepackage{amssymb}
\usepackage{textcomp}
\usepackage{graphicx}
\usepackage{booktabs}
\usepackage{tabularx}

\hypersetup{colorlinks=true,linkcolor=blue!60!black,urlcolor=teal}

% ── Farben ──────────────────────────────────────────────────
\definecolor{MerksatzFarbe}{HTML}{DDEEFF}
\definecolor{UebungFarbe}{HTML}{FFF3CC}
\definecolor{PraxisFarbe}{HTML}{E8F5E9}
\definecolor{TippFarbe}{HTML}{FFF8E1}
\definecolor{RosFarbe}{HTML}{FCE4EC}
\definecolor{AnalogieFarbe}{HTML}{EDE7F6}
\definecolor{LoesungFarbe}{HTML}{F3E5F5}
\definecolor{CodeHintergrund}{HTML}{F8F8F8}

% ── Python-Code ─────────────────────────────────────────────
\lstdefinestyle{pythonstyle}{
    language=Python,
    backgroundcolor=\color{CodeHintergrund},
    basicstyle=\ttfamily\small,
    keywordstyle=\color{blue!70!black}\bfseries,
    stringstyle=\color{teal!80!black},
    commentstyle=\color{gray!70!black}\itshape,
    showstringspaces=false,
    breaklines=true,
    frame=single,
    rulecolor=\color{gray!40},
    tabsize=4,
    extendedchars=true,
    numbers=left,
    numberstyle=\tiny\color{gray},
    numbersep=8pt
}
\lstset{
    style=pythonstyle,
    literate=
        {ä}{{\"a}}1 {ö}{{\"o}}1 {ü}{{\"u}}1
        {Ä}{{\"A}}1 {Ö}{{\"O}}1 {Ü}{{\"U}}1
        {ß}{{\ss}}1
}

% ── Didaktische Boxen ────────────────────────────────────────
\newtcolorbox{merksatz}[1][]{enhanced,breakable,
    colback=MerksatzFarbe,colframe=blue!60!black,
    fonttitle=\bfseries,title=Merksatz,#1}

\newtcolorbox{uebung}[1][]{enhanced,breakable,
    colback=UebungFarbe,colframe=orange!80!black,
    fonttitle=\bfseries,title=Übung,#1}

\newtcolorbox{praxis}[1][]{enhanced,breakable,
    colback=PraxisFarbe,colframe=green!60!black,
    fonttitle=\bfseries,title=Praxisbezug Robotik,#1}

\newtcolorbox{tipp}[1][]{enhanced,breakable,
    colback=TippFarbe,colframe=orange!40!black,
    fonttitle=\bfseries,title=Tipp,#1}

\newtcolorbox{ros}[1][]{enhanced,breakable,
    colback=RosFarbe,colframe=red!70!black,
    fonttitle=\bfseries,title=ROS2-Verbindung,#1}

\newtcolorbox{analogie}[1][]{enhanced,breakable,
    colback=AnalogieFarbe,colframe=purple!60!black,
    fonttitle=\bfseries,title=Analogie,#1}

\newtcolorbox{loesung}[1][]{enhanced,breakable,
    colback=LoesungFarbe,colframe=purple!60!black,
    fonttitle=\bfseries,title=Lösung,#1}
 in band*/main.tex
% ============================================================

% ── Buchreihe-Variablen (für shared/titelseite.tex kompatibel) ────
────────────────────────────────────────
\newcommand{\BuchReihe}{Robotik entwickeln lernen}
\newcommand{\BandAutor}{Wolfgang Becker}
\newcommand{\BandVersion}{Arbeitsfassung 2026}


% ── Sprache & Schrift ────────────────────────────────────────
\usepackage[ngerman]{babel}
\usepackage[T1]{fontenc}
\usepackage{lmodern}

% ── Layout ──────────────────────────────────────────────────
\usepackage{geometry}
\geometry{left=28mm,right=22mm,top=25mm,bottom=30mm}
\usepackage{scrlayer-scrpage}
\KOMAoptions{headsepline=true}
\pagestyle{scrheadings}
\clearpairofpagestyles
\lehead{\leftmark}
\rohead{\BuchReihe}   % wird in band*/main.tex gesetzt
\cfoot{\thepage}

% ── Pakete ──────────────────────────────────────────────────
\usepackage{xcolor}
\usepackage{tcolorbox}
\tcbuselibrary{skins,breakable,listings}
\usepackage{enumitem}
\usepackage{hyperref}
\usepackage{listings}
\usepackage{amsmath}
\usepackage{amssymb}
\usepackage{textcomp}
\usepackage{graphicx}
\usepackage{booktabs}
\usepackage{tabularx}

\hypersetup{colorlinks=true,linkcolor=blue!60!black,urlcolor=teal}

% ── Farben ──────────────────────────────────────────────────
\definecolor{MerksatzFarbe}{HTML}{DDEEFF}
\definecolor{UebungFarbe}{HTML}{FFF3CC}
\definecolor{PraxisFarbe}{HTML}{E8F5E9}
\definecolor{TippFarbe}{HTML}{FFF8E1}
\definecolor{RosFarbe}{HTML}{FCE4EC}
\definecolor{AnalogieFarbe}{HTML}{EDE7F6}
\definecolor{LoesungFarbe}{HTML}{F3E5F5}
\definecolor{CodeHintergrund}{HTML}{F8F8F8}

% ── Python-Code ─────────────────────────────────────────────
\lstdefinestyle{pythonstyle}{
    language=Python,
    backgroundcolor=\color{CodeHintergrund},
    basicstyle=\ttfamily\small,
    keywordstyle=\color{blue!70!black}\bfseries,
    stringstyle=\color{teal!80!black},
    commentstyle=\color{gray!70!black}\itshape,
    showstringspaces=false,
    breaklines=true,
    frame=single,
    rulecolor=\color{gray!40},
    tabsize=4,
    extendedchars=true,
    numbers=left,
    numberstyle=\tiny\color{gray},
    numbersep=8pt
}
\lstset{
    style=pythonstyle,
    literate=
        {ä}{{\"a}}1 {ö}{{\"o}}1 {ü}{{\"u}}1
        {Ä}{{\"A}}1 {Ö}{{\"O}}1 {Ü}{{\"U}}1
        {ß}{{\ss}}1
}

% ── Didaktische Boxen ────────────────────────────────────────
\newtcolorbox{merksatz}[1][]{enhanced,breakable,
    colback=MerksatzFarbe,colframe=blue!60!black,
    fonttitle=\bfseries,title=Merksatz,#1}

\newtcolorbox{uebung}[1][]{enhanced,breakable,
    colback=UebungFarbe,colframe=orange!80!black,
    fonttitle=\bfseries,title=Übung,#1}

\newtcolorbox{praxis}[1][]{enhanced,breakable,
    colback=PraxisFarbe,colframe=green!60!black,
    fonttitle=\bfseries,title=Praxisbezug Robotik,#1}

\newtcolorbox{tipp}[1][]{enhanced,breakable,
    colback=TippFarbe,colframe=orange!40!black,
    fonttitle=\bfseries,title=Tipp,#1}

\newtcolorbox{ros}[1][]{enhanced,breakable,
    colback=RosFarbe,colframe=red!70!black,
    fonttitle=\bfseries,title=ROS2-Verbindung,#1}

\newtcolorbox{analogie}[1][]{enhanced,breakable,
    colback=AnalogieFarbe,colframe=purple!60!black,
    fonttitle=\bfseries,title=Analogie,#1}

\newtcolorbox{loesung}[1][]{enhanced,breakable,
    colback=LoesungFarbe,colframe=purple!60!black,
    fonttitle=\bfseries,title=Lösung,#1}
